% ============================================================
% ch03.tex —— 第 3 章 素数：整数的原子
% ============================================================
\chapter{素数：整数的原子}

化学老师告诉你：世界上的物质千奇百怪，拆到底却只有一百多种原子；氢和氧拼成水，钠和氯拼成盐。整数世界有完全相同的故事：数有千千万，拆到底，都由同一批"基本粒子"拼成。这一章我们给整数做一次"光谱分析"——顺便证明，整数的"化学"比真实世界的化学还讲规矩：\textbf{同一种物质，配方全宇宙唯一}。

\begin{kaiwei}
把 $12$ 拆成更小因数的乘积，随便你怎么拆：
\[
12 = 3\times 4 = 3\times 2\times 2, \qquad
12 = 2\times 6 = 2\times 2\times 3 .
\]
两条路，终点站都是"一个 $3$、两个 $2$"。是巧合，还是定律？换 $60$ 试试：$60 = 6\times 10 = 2\times3\times2\times5$；换条路 $60 = 4 \times 15 = 2\times2\times3\times5$——又是同一个终点。
\end{kaiwei}

\section{拆不动的数}

先分类。有些数拆不动：$2, 3, 5, 7, 11, 13, 17, 19, 23, \dots$ 它们大于 $1$ 的因数只有自己。比如 $13$：试遍 $2$ 到 $12$，没有一个除得尽。这些"拆不动"的数是整数的\Keyword{原子}——正式名字：\textbf{素数}（prime number）。拆得动的（$4, 6, 8, 9, 10, 12, 14, 15, \dots$）叫\textbf{合数}，它们是"分子"，由原子拼成。素数清单开头是：
\[
2,\ 3,\ 5,\ 7,\ 11,\ 13,\ 17,\ 19,\ 23,\ 29,\ 31,\ 37,\ 41,\ 43,\ 47,\ 53,\ \cdots
\]
（$2$ 是唯一的偶素数——其余偶数都被 $2$ 整除，天生是合数。）

\textbf{$1$ 的身份最微妙}：它既拆不动（没有更小的因数），也当不了别人的零件。为什么干脆不开除它、也不选它当原子？看如果让 $1$ 加入素数会发生什么：
\[
6 = 2\times 3 = 1\times 2\times 3 = 1\times 1\times 2\times 3 = \cdots
\]
拆法突然变得有无穷多种，"配方唯一"瞬间报废。所以我们\textbf{发明规定}：$1$ 既不是素数，也不是合数。\textbf{定义是被设计的}：为了让好定理活得漂亮，我们修剪定义的边界——这是全书方法论的第一块展品（记住它，第 10 章"零为什么没有倒数"、第 32 章"椭圆曲线为什么要无重根"，全是同一个思路的重播）。

\section{大定理：拆法唯一}

\begin{dingli}[算术基本定理]
任何大于 $1$ 的整数，拆成素数的乘积，\textbf{不计次序，拆法唯一}。
\end{dingli}

比如 $360 = 2^3 \times 3^2 \times 5$，全世界的数学书都只能拆出这一种配方——"三个 $2$、两个 $3$、一个 $5$"。这个事实你从小用到大（约分、通分、找公因数），却可能从没想过它\textbf{需要证明}。"能拆"这一半好证；难的是"唯一"。我们分两步走。

\textbf{第一步：一个关键引理（素数的"穿透力"）。}

\begin{yinli}[欧几里得引理]
素数 $p$ 整除乘积 $ab$，则 $p$ 整除 $a$，或者 $p$ 整除 $b$（至少一个）。
\end{yinli}

先建立手感。$5$ 整除 $20 = 4\times 5$：确实整除了后一个因数。$7$ 整除 $42 = 6\times 7$：同样。$3$ 整除 $60 = 4 \times 15$：整除了 $15$。素数像一颗子弹，打进乘积里，\textbf{必打中其中一颗因子}——它不会"分摊火力"。合数就没这本事：$6$ 整除 $4\times 9 = 36$（$36 = 6\times 6$），但 $6$ 既不整除 $4$ 也不整除 $9$——它的火力拆成 "$2$ 打 $4$、$3$ 打 $9$" 两路分摊了。\textbf{穿透力是素数的独门绝技}。

现在证明。设 $p$ 是素数，$p \mid ab$。若 $p \mid a$，结论成立，收工。剩下要处理的情况是：$p$ 不整除 $a$。这时 $\gcd(p, a) = 1$——理由：$p$ 的正因数只有 $1$ 和 $p$，而 $p$ 不是 $a$ 的因数，所以两者的公因数只剩 $1$，即\textbf{互素}。第 2 章的裴蜀等式立刻上场：存在整数 $x, y$ 使
\[
xp + ya = 1 .
\]
（裴蜀系数要靠辗转相除倒代出来，第 2 章你练过：$4x + 5y = 1$ 的一组解是 $x = -1, y = 1$，即 $-4 + 5 = 1$。）

两边同乘 $b$：
\[
xpb + yab = b .
\]
（举个具体例子找感觉：$a = 4$，$p = 5$ 时，$4\times4 - 3\times5 = 16 - 15 = 1$，取 $x = -3$，$y = 4$——裴蜀系数靠第 2 章的"倒代"手艺造出来，通常一正一负，这里 $x = -3$ 是负的完全合法。）
看右边两项：第一项自带因子 $p$；第二项 $yab$ 中 $p \mid ab$（已知），所以第二项也是 $p$ 的倍数。两项之和——也就是 $b$——被 $p$ 整除。穿透力证明完毕。\textbf{第 2 章铺地砖的功夫，在这里第一次露出獠牙}：裴蜀等式不是绕口令，它是把"互素"变成可用等式的翻译器。

（顺带说透"分摊"为什么失败：$6 \mid 36$ 时，对 $a=4, b=9$ 用同样的流程会在"$\gcd(6,4)=2 \ne 1$"这一步卡死——$6$ 和 $4$ 不互素，裴蜀给不出 $6x + 4y = 1$，最多给出 $6x + 4y = 2$。穿透力依赖互素，素数的"因数只有 $1$ 和自己"保证了"不整除就互素"——这就是素数独门的原因。）

\textbf{第二步：唯一的证明。}假设 $n > 1$ 有两种素拆法：
\[
n = p_1 p_2 \cdots p_s = q_1 q_2 \cdots q_t
\]
（两边都是素数之积，可能重复，比如 $12 = 2\times2\times3$）。盯住 $p_1$：它整除右边的整体乘积 $q_1(q_2\cdots q_t)$。用引理：要么 $p_1 \mid q_1$，要么 $p_1 \mid q_2\cdots q_t$。若前者，$q_1$ 是素数，因数只有 $1$ 和自己，而 $p_1 > 1$，所以 $p_1 = q_1$；若后者，把引理再用一次（对更短的乘积），一路问下去，$p_1$ 必须等于某个 $q_j$。总之：\textbf{$p_1$ 必在右边名单里}，设它就是 $q_j$。两边同时划掉这个相同的素数，剩下
\[
p_2\cdots p_s = \text{（划掉 } q_j \text{ 的右边）}，
\]
两边仍然相等，且仍是素数之积。如法炮制：$p_2$ 也必在右边名单里，划掉；$p_3$……每一步配对消去一个。左边有限个，右边也有限个，最后必然\textbf{同时用完}（若一边先用完另一边还剩，就会出现"一堆素数之积等于 $1$"的荒谬，每个素数都 $\ge 2$，不可能）。所以两种拆法逐个配对、完全相同，只是排列次序不同。证毕。

\begin{cidian}
土名"原子"$\to$\textbf{素数}（prime）；土名"分子"$\to$\textbf{合数}；土名"配方"$\to$\textbf{素因数分解}（prime factorization）；这条大定理的两个名字——\textbf{算术基本定理}（Fundamental Theorem of Arithmetic）或\textbf{唯一分解定理}——说的是同一件事：分解存在，且唯一。
\end{cidian}

\section{配方在手，因数不愁}

唯一分解立刻升级你的因数技能：\textbf{不逐个试验，直接"数"出因数}。看 $360 = 2^3\times 3^2\times 5$。任何一个因数 $d$，拆开后也只能用这三种原子（否则 $d$ 里的新原子从哪来？），所以 $d$ 必形如 $2^a 3^b 5^c$，其中
\[
a \in \{0,1,2,3\}\ \text{（4 种）},\qquad b \in \{0,1,2\}\ \text{（3 种）},\qquad c \in \{0,1\}\ \text{（2 种）}.
\]
三种独立的小选择，乘法原理（穿三件衣服：$4\times3\times2$ 种搭配）：
\[
\text{因数总数} = (3+1)(2+1)(1+1) = 24 \text{ 个}.
\]
不用列一个验一个，一眼报出总数。\textbf{配方就是因数的户口本}。

再进一步把 $24$ 个因数全部"生成"出来——展开下面这个乘积，每一项恰好是一个因数：
\begin{align*}
(1+2+4+8)(1+3+9)(1+5) = {}& 1+2+3+4+5+6+8+9+10 \\
& +12+15+18+20+24+30+36+40+45 \\
& +60+72+90+120+180+360 .
\end{align*}
数一数，恰好 $24$ 项。\textbf{因数之和}也顺手拿下：$15\times13\times6 = 1170$。这套"生成器"技术（数学家叫它积性）是第 22 章欧拉乘积的直系祖先——你在初中做的因数展开，两千年后会长成素数定理的地基，这话不是修辞，是第 22 章的剧情预告。

\section{怎么判断一个大数是不是素数}

日常武器：\textbf{试除法}。判断 $n$ 是否素数，用不超过 $\sqrt{n}$ 的素数逐个去除。为什么只需试到平方根？因为因数成对出现：若 $n = ab$ 且 $a \le b$，则 $a\times a \le ab = n$，即 $a \le \sqrt{n}$——\textbf{大因数必然拖家带口地带一个小因数}，小因数藏在平方根以下。找遍平方根以下没找到搭档，就说明 $n$ 是素数。

实战：$91$ 是素数吗？$\sqrt{91} \approx 9.5$，只需试 $2, 3, 5, 7$：$91$ 是奇数（$\times$ 2）；$9+1 = 10$ 不是 $3$ 的倍数（$\times$ 3）；个位不是 $0/5$（$\times$ 5）；$91 = 13\times 7$（\checkmark 除尽！）——所以 $91 = 7\times 13$，合数。很多人凭"长得像素数"误判它。再战 $143$：$\sqrt{143}\approx 12$，试 $2,3,5,7,11$：$143 = 11\times 13$，合数，又是一个"伪装者"。这两个数是出题人的心头好，建议记住。

\begin{lishi}
欧几里得《几何原本》第七至九卷讨论素数，全书没有符号、没有方程，全部用文字和几何语言完成——包括"素数有无穷多"的证明（下一章主角）。值得一提的是"唯一分解"在《原本》里只以引理形式出现（第七卷命题 30 是欧几里得引理，配对的唯一性散落在几个命题里）；把它明确立为"基本定理"并命名，是高斯 1801 年《算术研究》的事。数学定理像建筑，奠基者和命名者常常不是同一人。
\end{lishi}

\begin{dongshou}
\begin{itemize}
  \yisi 分解 $2026$ 和 $1998$（用试除法，只试到平方根：$\sqrt{2026} \approx 45$，$\sqrt{1998}\approx 44.7$），验证第 2 章的 $\gcd(2026,1998) = 2$：两个配方的"公共部分"（每个素数取较低次幂）正是 $\gcd$。再顺手算 $\mathrm{lcm}(2026, 1998) = 2026\times 1998/2$（每个素数取较高次幂）。
  \yier 判断素数：$221$？$289$？$391$？（都试到平方根：$\sqrt{221}\approx 14.9$ 试 $2,3,5,7,11,13$；$289 = 17^2$，注意 $\sqrt{289} = 17$ 恰好要试到自己；$391 = 17\times 23$。）
  \yier 数一数 $720$ 有多少个因数，并求全部因数之和。（$720 = 2^4\cdot3^2\cdot5$；和 $= (1+2+4+8+16)(1+3+9)(1+5) = 31\times13\times6 = 2418$。用 3.3 节的生成器展开验算几个。）
  \yier 用"生成器"解释一个初中常识：完全平方数的因数个数是\textbf{奇数}，非完全平方数是偶数。（提示：因数配对 $d \leftrightarrow n/d$，谁和自己配对？对应到配方：$n$ 的每个素数指数都是偶数 $\iff$ 配对时有一项落单。）
\end{itemize}
\end{dongshou}

\begin{shaonao}
两个数的 $\gcd$ 与最小公倍数 $\mathrm{lcm}$ 之间有个漂亮的等式：
\[
\gcd(a,b) \times \mathrm{lcm}(a,b) = a \times b .
\]
先用 $12$ 和 $18$ 验证（$6\times36 = 216 = 12\times18$ \checkmark）。然后用素数配方证明它：把两个配方并排写
\[
a = p_1^{e_1}\cdots p_k^{e_k},\qquad b = p_1^{f_1}\cdots p_k^{f_k}
\]
（指数允许为 $0$，把两边补齐到同一批素数），$\gcd$ 取每个素数的 $\min(e_i, f_i)$，$\mathrm{lcm}$ 取 $\max(e_i, f_i)$。关键一步是发现恒等式 $\min(e,f) + \max(e,f) = e + f$（一个取小一个取大，合起来正好是两个的和）——整个定理就压缩在这一行里。证完你会发现：$\gcd\times\mathrm{lcm} = ab$ 不是巧合，是"$\min$ 加 $\max$ 等于和"穿了件数学外套。
\end{shaonao}

\begin{chengshi}
唯一分解的严格证明用"最小反例法"或数学归纳法把第二步的"如法炮制"说严谨（严格说要对 $n$ 归纳：消去一个素数后剩下的乘积更小，归纳假设适用），本书从略；欧几里得引理的证明我们是完整给出的——它是全章的技术核心。"任何大于 $1$ 的整数都能拆到只剩素数"依赖一个朴素事实："合数必有小于自身的真因数"，反复拆数越来越小，不能永远拆下去（正整数的降链必止）。
\end{chengshi}

下一章问一个孩子都会问、两千三百年前就被认真回答过的问题：素数，到底有没有用完的一天？答案的证明只有五行，却被称为数学史上最漂亮的证明之一。
